\subsection{Ray tracing examples}

This first example is intended to be a first contact OASYS2 and SHADOW4. Just experience using widgets and making simple simulations. 

\textbf{You will:}

\begin{itemize}
    \item define grid and geometrical sources
    \item Experience with different mirrors
    \item find the position of foci
    \item study the image (intensity distribution) at image and on the mirror (footprint)
\end{itemize}

\textbf{Questions:}

\textbf{i)}  

Define grid source over a rectangle in real space and collimated in divergence space

Experience with different mirrors to focus it and find the focal positions. The following mirrors could be studied: 

\begin{itemize}
    \item Plane mirror (no focus)
    \item Spherical mirror in normal incidence (stigmatic focus)
    \item Spherical mirror in not-normal (grazing) incidence (astigmatic focus)
    \item Parabolic mirror, not-normal incidence (perfect stigmatic focus)
\end{itemize}

\textbf{ii)}  

\begin{itemize}
    \item define geometrical source, point in real space, and isotropic (cone) in divergence space
    \item place a spherical mirror of $R$=1 m and distances $p=q=$ 1 m. See perfect focus and homogeneous footprint.
\end{itemize}


\textbf{Answer}

Use the OASYS workspace file {\tt introduction.ows}. Define grid with the center sligtly shifted from (0,0). Define mirrors with internal calculation of the mirror parameters using a focal distance $q=$0.25 m. Use ``FocNew" widget to see the evolution of the beam size, and to check is the focus is at the image position $q$. Results for Spherical mirror with 45$^\circ$ incidence are at Fig. 1. 

\begin{figure}[H]
\label{fig:sf-fig1}
\caption{Left: beam evolution (from FocNew). In vertical (Z), the best focus is at zero (at the image plane). In horizontal (X) the best focus is at 0.25 m downstream from the image plane, as a result of astigmatism provoked by the non-normal incidence with spherical reflectors. Right: diagram of the path of the rays.}
\includegraphics[width=0.48\textwidth]{shadow4-intro/fig1a.png}
\includegraphics[width=0.48\textwidth]{shadow4-intro/fig1b.png}
\end{figure}



